\section{Open questions}
\label{sec:discussion}

\begin{enumerate}\setlength\itemsep{2pt}
\item \emph{Hull versus point set.} Is $\operatorname{Iso}(\mathcal G(S))=\operatorname{Iso}(P_S)$ for every magic, or semi-magic, square $S$, or more generally for every arrangement whose heights have vanishing first moments? By \Cref{prop:pointset} this would determine the isometry groups of all these gems; it holds for order four (\Cref{prop:iso4}) and for the Siamese and digit-exchanged squares (\Cref{thm:siamese}, \Cref{prop:exchanged}). For odd $n$ the self-complementary squares are associative (\Cref{lem:selfcomp}), so the point set of an order-five square has isometry group $\{\pm\mathrm{id}\}$ or trivial; whether the same holds for the gems is open, and an exhaustive check over the $275{,}305{,}224$ classes of order five would take about seven CPU-weeks at the rate of \Cref{sec:census-4}.
\item \emph{Growth of the vertex count for arbitrary arrangements.} For uniformly random arrangements the mean vertex count grows from $8.22$ at $n=3$ to $27.91$ at $n=10$, and the ratio $\mathbb E[V]/\ln^2N$ decreases slowly from $1.70$ to $1.32$ (\Cref{tab:vertices}; estimates). For $N$ independent uniform points in a $3$-polytope the expected vertex count is $\Theta(\ln^2N)$ \cite{barany1993random}; our points are not independent (the grid is fixed and only the heights are permuted), so this is only a heuristic analogy. Is $\mathbb E[V]=\Theta(\ln^2 N)$ for random arrangements? By \Cref{sec:magic-vs-random} the line conditions hardly change the mean vertex count at orders four and five; do they change its growth rate?
\item \emph{Secondary cones.} By \Cref{prop:subdivisions} the face lattice of the gem of a height function $z$ is fixed by the secondary cones containing $z$ and $-z$, so the combinatorics of the gems of a family of arrangements is the set of pairs of chambers of the secondary fan of the grid that it meets. Which pairs are met by the permuted heights that satisfy the line conditions?
\item \emph{Linear squares.} The shear of \Cref{thm:siamese} reduces the lift of a linear square with fixed integer vectors $\mathbf u,\mathbf w$ to a number of levels that does not depend on $n$, but the hull of the lattice points of a level polygon may a priori have more and more vertices as $n$ grows. Is the vertex count of such a family bounded independently of $n$, as it is for the Siamese and the digit-exchanged squares? How does the largest vertex count of a linear magic square of order $n$ ($18$, $25$, $31$, $36$ for $n=5,7,9,11$; \Cref{rem:linear}) grow?
\end{enumerate}
